\section{Kết xuất thể tích}
\label{subsec:volume_rendering_detailed}

Kết xuất thể tích (Volume Rendering) là một tập hợp các kỹ thuật trong đồ họa máy tính được sử dụng để tạo ra hình ảnh 2D từ một bộ dữ liệu 3D được lấy mẫu rời rạc (thường là một lưới 3D vô hướng) hoặc từ một hàm mô tả thuộc tính của không gian 3D một cách liên tục. Khác với kỹ thuật kết xuất bề mặt truyền thống chỉ hiển thị hình học bề mặt của đối tượng, kết xuất thể tích có khả năng mô tả và hiển thị các hiện tượng không có bề mặt rõ ràng như mây, khói, sương mù, lửa, cũng như dữ liệu y tế mà không cần trích xuất bề mặt hình học một cách tường minh \cite{Kajiya84Ray}.

Nền tảng vật lý phổ biến nhất cho kết xuất thể tích dựa trên mô hình tương tác giữa ánh sáng và môi trường. Khi một tia sáng đi xuyên qua môi trường thể tích, cường độ và màu sắc của nó thay đổi do hai quá trình chính: sự phát xạ ánh sáng từ các hạt trong môi trường và sự suy giảm cường độ do hấp thụ. 

Xét một tia nhìn được tham số hóa bởi $r(s) = o + s\mathbf{v}$, trong đó $o$ là vị trí gốc (thường là tâm máy ảnh), $\mathbf{v}$ là vector đơn vị chỉ hướng tia, và $s$ là khoảng cách dọc theo tia. Giả sử tia đi vào môi trường tại $s = s_{near}$ và đi ra tại $s = s_{far}$.

Tại mỗi điểm dọc theo tia $r(s)$, môi trường được đặc trưng bởi hai thuộc tính chính:
\begin{itemize}
    \item Mật độ thể tích $\sigma(r(s))$: Một đại lượng vô hướng không âm ($\sigma \ge 0$) mô tả mức độ "dày đặc" của môi trường tại điểm đó. Nó đại diện cho xác suất vi phân mà một photon tương tác (bị hấp thụ) trên một đơn vị độ dài. Khi $\sigma = 0$, môi trường hoàn toàn trong suốt; khi $\sigma > 0$, môi trường có khả năng phát xạ và hấp thụ.
    \item Bức xạ phát xạ $\mathcal{L}_{emitted}(r(s), \mathbf{v})$: Đại diện cho lượng ánh sáng được phát ra từ điểm đó theo hướng nhìn $\mathbf{v}$.
\end{itemize}

Để mô tả sự suy giảm ánh sáng dọc theo tia, ta định nghĩa độ truyền qua tích lũy $\mathcal{T}(s_{near}, s)$ – mô tả xác suất ánh sáng có thể đi từ điểm $s_{near}$ đến $s$ mà không bị hấp thụ:
\begin{equation} \label{eq:transmittance_cont}
\mathcal{T}(s_{near}, s) = \exp\left(-\int_{s_{near}}^{s} \sigma(r(u)) \, du\right)
\end{equation}

Giá trị $\mathcal{T}$ giảm dần theo hàm mũ khi tia đi sâu vào môi trường có mật độ $\sigma > 0$.

Kết hợp các yếu tố trên, phương trình kết xuất thể tích mô tả tổng lượng ánh sáng (bức xạ) mà máy ảnh thu nhận được dọc theo một tia nhìn $r(s) = o + s\mathbf{v}$ đi qua một môi trường thể tích:

\begin{equation} \label{eq:volume_integral_detailed}
    \mathcal{L}_{observed}(r) = \int_{s_{near}}^{s_{far}} \underbrace{\mathcal{T}(s_{near}, s)}_{\text{(1) Độ truyền qua}} \cdot \underbrace{\sigma(r(s))}_{\text{(2) Mật độ thể tích}} \cdot \underbrace{\mathcal{L}_{emitted}(r(s), \mathbf{v})}_{\text{(3) Bức xạ phát xạ}} \, ds
\end{equation}